\documentclass[11pt]{amsart}

\usepackage[a4paper,margin=1in]{geometry}
\usepackage{amsmath,amssymb,amsthm,mathtools}
\usepackage{booktabs}
\usepackage{enumitem}
\usepackage{graphicx}
\usepackage[hidelinks]{hyperref}

\title{An Ordinal-Prime Transform of the M\"obius Function}
\author{}
\date{Version 3}

\newtheorem{theorem}{Theorem}[section]
\newtheorem{proposition}[theorem]{Proposition}
\newtheorem{lemma}[theorem]{Lemma}
\newtheorem{corollary}[theorem]{Corollary}
\theoremstyle{definition}
\newtheorem{definition}[theorem]{Definition}
\newtheorem{remark}[theorem]{Remark}
\newtheorem{heuristic}[theorem]{Heuristic}

\newcommand{\mob}{\mu}
\newcommand{\Mert}{M}
\newcommand{\Sq}{Q}
\newcommand{\PS}{\mathcal P}

\begin{document}

\maketitle
\tableofcontents

\section{Purpose and scope}

This note develops the ordinal-prime transform of the M\"obius function and
determines exactly how much arithmetic information it carries. It is
self-contained and uses no other construction.

The starting point is the M\"obius function $\mob(n)\in\{-1,0,+1\}$. The
occurrences of $-1$ and of $+1$ are numbered separately, in increasing order of
$n$, and the corresponding ordinal prime is attached to each occurrence; zeros
are mapped to $1$. The resulting signed sequence is $S(n)$ and its partial sum
is $SS(n)=\sum_{k\le n}S(k)$.

The conclusions are:

\begin{enumerate}[label=(\roman*)]
\item $SS(n)$ has an exact closed form in terms of the Mertens function
$\Mert(n)$, the squarefree counting function $\Sq(n)$, and the sum of the first
$m$ primes (Proposition~\ref{prop:closed}).
\item After subtracting the deterministic zero contribution $n-\Sq(n)$, the
remainder $T(n):=SS(n)-n+\Sq(n)$ is \emph{exactly} a signed sum of $|\Mert(n)|$
consecutive primes, occupying a window of prime indices symmetric about
$(\Sq(n)+1)/2$ (Proposition~\ref{prop:block}).
\item Consequently $T(n)=\Mert(n)\,p_{\lfloor \Sq(n)/2\rfloor}(1+\theta(n))$
with $\theta(n)\to0$ effectively, and the Mertens criterion transfers to $T$ in
both directions (Theorem~\ref{thm:RH}).
\item Every object built here --- $SS$, $T$, the residual $E$, and the centered
residual $E_{\mathrm{mid}}$ --- is a deterministic function of the pair
$(\Mert(n),\Sq(n))$ alone. The transform therefore creates no independent
arithmetic constraint on $\Mert(n)$. This applies to the residuals as well and
is stated as Proposition~\ref{prop:info}.
\item The nonlinearity of the transform splits exactly into a central part
linear in $\Mert(n)$ and an \emph{antisymmetric weighted sum of prime gaps}
with triangular weights (Theorem~\ref{thm:gaps}). The weight norm has a closed
form (Lemma~\ref{lem:w2}), which forces the exponent $3/2$ in the observed law
$|E_{\mathrm{mid}}(n)|\approx0.155\,|\Mert(n)|^{3/2}\log n$.
\item The size of that residual is governed by the local covariance structure of
prime gaps, not by the arithmetic of $\mob$. A model using the measured gap
covariance predicts its standard deviation to within $0.3\%$ with no free
parameter, whereas an independent-gap model underpredicts by a factor $0.853$
(Section~\ref{sec:numeric}, part D).
\end{enumerate}

All numerical statements in this note were verified for $1\le n\le10^{6}$; the
verification protocol and effective ranges are given in
Section~\ref{sec:numeric}.

\section{Notation and elementary lemmas}

Let $p_1=2,p_2=3,p_3=5,\ldots$ denote the primes in increasing order, and set
$p_0:=0$ by convention (used only to make the product forms below well defined
at $n=1$).

Write
\[
\Mert(n):=\sum_{k\le n}\mob(k),
\qquad
\Sq(n):=\sum_{k\le n}\mob(k)^2 ,
\qquad
\PS(m):=\sum_{j=1}^{m}p_j ,\quad \PS(0):=0 .
\]

Define the cumulative occurrence counts
\[
N_-(n)=\sum_{k\le n}\frac{\mob(k)^2-\mob(k)}2,
\qquad
N_0(n)=\sum_{k\le n}\bigl(1-\mob(k)^2\bigr),
\qquad
N_+(n)=\sum_{k\le n}\frac{\mob(k)^2+\mob(k)}2 .
\]
Thus $N_-(n)$, $N_0(n)$, $N_+(n)$ count the occurrences of $\mob=-1$, $\mob=0$,
$\mob=+1$ up to $n$, and the counts are inclusive of $n$ itself.

\begin{lemma}[Counter identities]\label{lem:counters}
For every $n\ge1$,
\[
N_+(n)+N_-(n)=\Sq(n),
\qquad
N_+(n)-N_-(n)=\Mert(n),
\qquad
N_0(n)=n-\Sq(n),
\]
hence
\[
N_+(n)=\frac{\Sq(n)+\Mert(n)}2,
\qquad
N_-(n)=\frac{\Sq(n)-\Mert(n)}2 .
\]
\end{lemma}

\begin{lemma}[Parity]\label{lem:parity}
$\Sq(n)\equiv \Mert(n)\pmod 2$ for every $n\ge1$. Consequently the indices
$(\Sq(n)\pm\Mert(n))/2$ occurring throughout are integers.
\end{lemma}

\begin{proof}
$\Sq=N_++N_-$ and $\Mert=N_+-N_-$, so $\Sq-\Mert=2N_-$.
\end{proof}

Lemma~\ref{lem:parity} is used silently in every subsequent formula and was
absent from version 1; without it the central identities are not well formed.

\section{The transform}

\begin{definition}
The three rows are
\[
R_-(n)=-\frac{\mob(n)^2-\mob(n)}2\,p_{N_-(n)},
\qquad
R_0(n)=1-\mob(n)^2,
\qquad
R_+(n)=\frac{\mob(n)^2+\mob(n)}2\,p_{N_+(n)} ,
\]
with the convention $p_0=0$ (needed only at $n=1$, where $N_-(1)=0$ and the
coefficient vanishes). The ordinal-prime transform is
\[
S(n):=R_-(n)+R_0(n)+R_+(n)
=
\begin{cases}
-p_{N_-(n)},&\mob(n)=-1,\\[1mm]
1,&\mob(n)=0,\\[1mm]
+p_{N_+(n)},&\mob(n)=+1 .
\end{cases}
\]
\end{definition}

Along the occurrences of $\mob=-1$ the values are $-2,-3,-5,-7,-11,\ldots$; along
the occurrences of $\mob=+1$ they are $2,3,5,7,11,\ldots$, in a separate ordinal
count. Explicitly,
\[
S(1),\ldots,S(12)=2,\,-2,\,-3,\,1,\,-5,\,3,\,-7,\,1,\,1,\,5,\,-11,\,1 .
\]

\begin{definition}
$SS(n):=\sum_{k=1}^{n}S(k)$ and $T(n):=SS(n)-\bigl(n-\Sq(n)\bigr)$.
\end{definition}

\begin{proposition}[Closed form]\label{prop:closed}
For every $n\ge1$,
\[
\boxed{
SS(n)=\PS\bigl(N_+(n)\bigr)-\PS\bigl(N_-(n)\bigr)+N_0(n)
=
\PS\!\left(\frac{\Sq(n)+\Mert(n)}2\right)
-
\PS\!\left(\frac{\Sq(n)-\Mert(n)}2\right)
+n-\Sq(n).
}
\]
\end{proposition}

\begin{proof}
The $k\le n$ with $\mob(k)=+1$ receive, in order, the primes
$p_1,\dots,p_{N_+(n)}$, contributing $\PS(N_+(n))$; those with $\mob(k)=-1$
contribute $-\PS(N_-(n))$; each of the $N_0(n)$ zeros contributes $1$. The
second form is Lemma~\ref{lem:counters}, the indices being integers by
Lemma~\ref{lem:parity}.
\end{proof}

\section{The signed prime block}

\begin{proposition}[Block representation]\label{prop:block}
Write $W(n)$ for the set of prime indices strictly between $N_-(n)$ and
$N_+(n)$, inclusive of the larger endpoint; $W(n)$ has exactly $|\Mert(n)|$
elements and is symmetric about the index $(\Sq(n)+1)/2$. Then
\[
T(n)=\operatorname{sgn}\bigl(\Mert(n)\bigr)\sum_{j\in W(n)}p_j
=
\begin{cases}
\displaystyle\sum_{j=(\Sq-\Mert)/2+1}^{(\Sq+\Mert)/2}p_j, & \Mert(n)>0,\\[5mm]
0,& \Mert(n)=0,\\[2mm]
\displaystyle-\sum_{j=(\Sq-|\Mert|)/2+1}^{(\Sq+|\Mert|)/2}p_j, & \Mert(n)<0 .
\end{cases}
\]
In particular
\[
\operatorname{sgn}T(n)=\operatorname{sgn}\Mert(n),
\qquad
T(n)=0\iff \Mert(n)=0 .
\]
\end{proposition}

\begin{proof}
Immediate from Proposition~\ref{prop:closed}: $T=\PS(N_+)-\PS(N_-)$ is a
telescoping difference of prefix sums of the primes. The window endpoints are
$\min(N_\pm)+1$ and $\max(N_\pm)$, whose midpoint is
$(N_++N_-+1)/2=(\Sq+1)/2$.
\end{proof}

\begin{corollary}[Exact bounds]\label{cor:bounds}
\[
|\Mert(n)|\,p_{\frac{\Sq(n)-|\Mert(n)|}{2}+1}
\;\le\;
|T(n)|
\;\le\;
|\Mert(n)|\,p_{\frac{\Sq(n)+|\Mert(n)|}{2}} .
\]
\end{corollary}

Both inequalities are trivial (each of the $|\Mert(n)|$ primes in the block lies
between the smallest and the largest one), but the \emph{lower} bound is what
makes the transfer of Mertens bounds an equivalence rather than an implication;
see Theorem~\ref{thm:RH}.

\section{The prime scale, with explicit error}

We record the multiplicative estimate in a form that is valid at every $n$,
including the infinitely many $n$ with $\Mert(n)=0$. (Version 1 wrote
$T(n)\sim \Mert(n)p_{\Sq(n)/2}$; that assertion is empty on the zero set of
$\Mert$, which contains $5361$ values of $n\le10^{6}$.)

\begin{proposition}\label{prop:scale}
Define $\theta(n):=0$ if $\Mert(n)=0$ and otherwise by
$T(n)=\Mert(n)\,p_{\lfloor \Sq(n)/2\rfloor}\bigl(1+\theta(n)\bigr)$. Then
\[
|\theta(n)|\;\le\;
\max_{|j-\Sq(n)/2|\le |\Mert(n)|/2+1}
\left|\frac{p_j}{p_{\lfloor \Sq(n)/2\rfloor}}-1\right| ,
\]
and $\theta(n)\to0$ as $n\to\infty$.
\end{proposition}

\begin{proof}
The bound is Proposition~\ref{prop:block} together with the triangle
inequality. For the limit: by the prime number theorem in the form
$\Mert(n)=o(n)$, and by $\Sq(n)=\tfrac{6}{\pi^2}n+O(\sqrt n)$, one has
$|\Mert(n)|/\Sq(n)\to0$; hence the window endpoints are
$(1+o(1))\Sq(n)/2$ on the index scale, and $p_{(1+o(1))m}/p_m\to1$ by the prime
number theorem again.
\end{proof}

Both hypotheses invoked here are named: $\Mert(n)=o(n)$ is equivalent to the
prime number theorem, and $\Sq(n)=\tfrac6{\pi^2}n+O(\sqrt n)$ is elementary.

\begin{remark}[The asymptotic constant is a poor numerical guide]
From $p_m=m(\log m+\log\log m-1+o(1))$ and $\Sq(n)/2\sim \tfrac3{\pi^2}n$ one
gets
\[
p_{\lfloor \Sq(n)/2\rfloor}
=\frac{3}{\pi^2}\,n\log n\left(1+\frac{\log\log n-1+\log(3/\pi^2)}{\log n}+o\!\left(\tfrac1{\log n}\right)\right),
\]
so the correction to the leading constant decays only like $1/\log n$. Measured
values of $T(n)\big/\bigl(\tfrac3{\pi^2}n\log n\,\Mert(n)\bigr)$ are
$0.9558,\ 0.9954,\ 1.0151,\ 1.0281$ at $n=10^3,10^4,10^5,10^6$ --- not monotone
and still $2.8\%$ off at $n=10^6$. By contrast
$T(n)/\bigl(\Mert(n)p_{\lfloor \Sq(n)/2\rfloor}\bigr)$ equals
$1.00200,\ 1.00071,\ 1.000097,\ 0.999998$ at the same points. Any quantitative
use should keep $p_{\lfloor \Sq(n)/2\rfloor}$ and not the closed constant.
\end{remark}

\section{Transfer of the Mertens criterion}

\begin{proposition}
$|T(n)|=O\bigl(n\log n\,|\Mert(n)|\bigr)$ and hence
$SS(n)=O\bigl(n\log n\,|\Mert(n)|+n\bigr)$. Unconditionally
$SS(n)=o(n^2\log n)$, by $\Mert(n)=o(n)$ (prime number theorem).
\end{proposition}

\begin{theorem}\label{thm:RH}
The following are equivalent.
\begin{enumerate}[label=(\alph*)]
\item The Riemann hypothesis.
\item $\Mert(n)=O_\varepsilon(n^{1/2+\varepsilon})$ for every $\varepsilon>0$.
\item $SS(n)-n+\Sq(n)=O_\varepsilon(n^{3/2+\varepsilon})$ for every
$\varepsilon>0$.
\end{enumerate}
\end{theorem}

\begin{proof}
(a)$\iff$(b) is classical (Littlewood).

(b)$\Rightarrow$(c): by Corollary~\ref{cor:bounds},
$|T(n)|\le|\Mert(n)|p_{(\Sq+|\Mert|)/2}\ll |\Mert(n)|\,n\log n
\ll_\varepsilon n^{3/2+\varepsilon}\log n$, and the logarithm is absorbed into
$n^{\varepsilon}$.

(c)$\Rightarrow$(b): this is the direction version 1 asserted without proof. By
the prime number theorem $|\Mert(n)|=o(n)$ and $\Sq(n)=\tfrac6{\pi^2}n+O(\sqrt
n)$, so
\[
\frac{\Sq(n)-|\Mert(n)|}{2}+1=\bigl(1+o(1)\bigr)\frac{3}{\pi^2}n ,
\]
whence the smallest prime of the block satisfies
$p_{(\Sq-|\Mert|)/2+1}\ge c\,n\log n$ for some $c>0$ and all large $n$. The
lower bound of Corollary~\ref{cor:bounds} then gives
\[
|\Mert(n)|\le \frac{|T(n)|}{c\,n\log n}
\ll_\varepsilon \frac{n^{3/2+\varepsilon}}{n\log n}
\ll_\varepsilon n^{1/2+\varepsilon}. \qedhere
\]
\end{proof}

\begin{remark}
Theorem~\ref{thm:RH} is a reformulation, not a proof mechanism. Note also that
(b) must not be strengthened to $\Mert(n)=O(\sqrt n)$: the Riemann hypothesis
does not imply the Mertens conjecture, which is false (Odlyzko--te Riele, 1985).
\end{remark}

\section{The zeta ingredients}

The two M\"obius-derived ingredients have the classical Dirichlet series
\[
\sum_{m\ge1}\frac{\mob(m)}{m^s}=\frac1{\zeta(s)},
\qquad
\sum_{m\ge1}\frac{\mob(m)^2}{m^s}=\frac{\zeta(s)}{\zeta(2s)},
\qquad \Re s>1,
\]
and the prime weights are governed by $\log\zeta(s)$ through the prime-counting
function, via the exact Abel identity
\[
\PS(m)=m\,p_m-\int_2^{p_m}\pi(t)\,dt ,
\]
(verified here for $m\le100$; e.g.\ $\PS(3)=10=3\cdot5-5$).

Version 1 substituted Perron integrals for $\Mert$ and $\Sq$ into the arguments
of $\PS$. That is not a functional identity and it is dropped here, for two
reasons. First, $\PS$ is defined only on nonnegative integers, so an integral
placed inside its argument is notation for ``the integer that this integral
evaluates to'', not an analytic representation. Second, at integer $x$ the
untruncated Perron integral returns $\Mert(x)-\tfrac12\mob(x)$, not $\Mert(x)$.

The honest statement is the following. The transform composes three classical
objects,
\[
\frac1{\zeta(s)}\ \ (\text{signed M\"obius imbalance}),
\qquad
\frac{\zeta(s)}{\zeta(2s)}\ \ (\text{squarefree count}),
\qquad
\log\zeta(s)\ \ (\text{prime weights}),
\]
through the nonlinear expression
$\PS\bigl(\tfrac{\Sq+\Mert}2\bigr)-\PS\bigl(\tfrac{\Sq-\Mert}2\bigr)+n-\Sq$.
None of the three components is new, and the composition is not a Dirichlet
series of anything.

\section{Exact information content}

\begin{proposition}\label{prop:info}
For each $n$, $SS(n)$ is a deterministic function of the pair
$(\Mert(n),\Sq(n))$; and with $\Sq(n)$ held fixed, $T(n)$ is strictly
increasing in $\Mert(n)$. The same is true of every residual defined below,
including $E(n)$ and $E_{\mathrm{mid}}(n)$.
\end{proposition}

\begin{proof}
The first claim is Lemma~\ref{lem:counters} plus
Proposition~\ref{prop:closed}. For the second, increasing $\Mert$ by $2$ at
fixed $\Sq$ raises $N_+$ by one and lowers $N_-$ by one, changing $T$ by
$p_{N_++1}+p_{N_-}>0$. The residuals $E$ and $E_{\mathrm{mid}}$ are $T$ minus an
explicit function of $(\Mert,\Sq)$, hence also functions of $(\Mert,\Sq)$.
\end{proof}

Two consequences deserve to be stated plainly.

First, $SS(n)$ does not retain the internal ordering of the $\pm1$ values of
$\mob$ below $n$: two M\"obius sign histories with the same endpoint counts
produce the same $SS(n)$. It is a nonlinear endpoint statistic.

Second --- and this corrects a claim in version 1 --- the residual $E(n)$ is
\emph{not} structurally distinct from $SS(n)$ in the informational sense. Its
window is determined by $(\Mert(n),\Sq(n))$, so the argument above applies to it
verbatim. What is genuinely different about the residual is not its information
content but its \emph{size}, which turns out to be governed by prime-gap
statistics; that is the subject of the next two sections.

\section{The residual and its correct centering}

\begin{definition}
Let
\[
p^{*}(n):=\tfrac12\Bigl(p_{\lfloor (\Sq(n)+1)/2\rfloor}+p_{\lceil (\Sq(n)+1)/2\rceil}\Bigr)
\]
be the linear interpolation of the prime sequence at the exact window centre
$(\Sq(n)+1)/2$ of Proposition~\ref{prop:block}. Define
\[
E(n):=T(n)-\Mert(n)\,p_{\lfloor \Sq(n)/2\rfloor},
\qquad
E_{\mathrm{mid}}(n):=T(n)-\Mert(n)\,p^{*}(n),
\]
so that the exact decomposition
\[
\boxed{\;T(n)=\Mert(n)\,p^{*}(n)+E_{\mathrm{mid}}(n)\;}
\]
holds for every $n$.
\end{definition}

When $\Sq(n)$ is odd, $(\Sq+1)/2$ is an integer and $p^{*}=p_{(\Sq+1)/2}$; when
$\Sq(n)$ is even, $p^{*}$ is the mean of the two central primes.

\begin{proposition}[Centering offset]\label{prop:offset}
$E(n)-E_{\mathrm{mid}}(n)=\Mert(n)\bigl(p^{*}(n)-p_{\lfloor \Sq(n)/2\rfloor}\bigr)$,
where the bracket equals one half of a prime gap when $\Sq(n)$ is even and one
full prime gap when $\Sq(n)$ is odd. It is positive and of size
$\asymp\log p^{*}(n)$, so $E$ carries a deterministic component linear in
$\Mert(n)$ which $E_{\mathrm{mid}}$ does not.
\end{proposition}

The identity was verified exactly for all $n\le10^{6}$; the median of
$\bigl(p^{*}-p_{\lfloor \Sq/2\rfloor}\bigr)/\log p_{\lfloor \Sq/2\rfloor}$ over
$10^{4}\le n\le10^{6}$ is $0.5318$. Proposition~\ref{prop:offset} explains the
correlation structure of the two residuals: the offset is a fixed positive
multiple of $\Mert(n)$, so $E$ correlates with $\Mert$ and $E_{\mathrm{mid}}$
does not. Measured over $100\le n\le10^{6}$,
\[
\operatorname{corr}(E,\Mert)=+0.2573,
\qquad
\operatorname{corr}(E_{\mathrm{mid}},\Mert)=+0.0045 .
\]
$E_{\mathrm{mid}}$ is therefore the object to study; $E$ is $E_{\mathrm{mid}}$
plus a bookkeeping artefact. Version 1 studied only $E$.

\section{Exact antisymmetric gap representation}

Write $g_i:=p_{i+1}-p_i$ and $m(n):=\lfloor(\Sq(n)+1)/2\rfloor$.

\begin{theorem}[Gap representation]\label{thm:gaps}
Let $h:=|\Mert(n)|$. Then
\[
\boxed{\;
E_{\mathrm{mid}}(n)=\operatorname{sgn}\bigl(\Mert(n)\bigr)\sum_{i}w_i\,g_i \;}
\]
where the weights $w_i$ are supported on $h-1$ or $h-2$ gaps symmetrically
placed about $m=m(n)$, are antisymmetric in sign about the centre, and are given
explicitly by
\[
h \text{ even}:\quad
w_{m+k}=\tfrac h2-k,\quad w_{m-k}=-\bigl(\tfrac h2-k\bigr),\quad 1\le k\le \tfrac h2-1;
\]
\[
h \text{ odd}:\quad
w_{m+k}=\tfrac{h-1}2-k \ \ (0\le k\le \tfrac{h-3}2),
\qquad
w_{m-k}=-\bigl(\tfrac{h+1}2-k\bigr)\ \ (1\le k\le \tfrac{h-1}2).
\]
In both cases the gap at the centre cancels identically, and the multiset of
weight magnitudes on each side is $\{1,2,\dots,\lfloor (h-1)/2\rfloor\}$.
\end{theorem}

\begin{proof}
Take $\Mert(n)>0$; the other sign is identical after a global change of sign.

Suppose $h$ is even, so $\Sq$ is even by Lemma~\ref{lem:parity}, $m=\Sq/2$ and
$p^{*}=\tfrac12(p_m+p_{m+1})$. The window is $\{m-\tfrac h2+1,\dots,m+\tfrac h2\}$.
For $j>m$, $p_j-p^{*}=\tfrac12 g_m+\sum_{i=m+1}^{j-1}g_i$; for $j\le m$,
$p_j-p^{*}=-\tfrac12 g_m-\sum_{i=j}^{m-1}g_i$. Summing over the window, the
$\tfrac12 g_m$ terms occur $\tfrac h2$ times with each sign and cancel. Counting
how many $j$ contain a given $g_{m+k}$ gives the coefficient $\tfrac h2-k$, and
symmetrically $-(\tfrac h2-k)$ for $g_{m-k}$.

If $h$ is odd then $\Sq$ is odd, $m=(\Sq+1)/2$ and $p^{*}=p_m$, the window being
$\{m-\tfrac{h-1}2,\dots,m+\tfrac{h-1}2\}$. The same count gives the stated
coefficients, the two sides again carrying magnitudes
$\tfrac{h-1}2,\dots,1$.
\end{proof}

\begin{lemma}[Weight norm]\label{lem:w2}
\[
\boxed{\;
\sum_i w_i^{2}=
\begin{cases}
\dfrac{h(h^{2}-1)}{12}, & h \text{ odd},\\[3mm]
\dfrac{h(h-1)(h-2)}{12}, & h \text{ even},
\end{cases}
\qquad\text{so}\qquad
\sum_i w_i^{2}=\frac{h^{3}}{12}\bigl(1+O(h^{-1})\bigr). }
\]
\end{lemma}

\begin{proof}
By Theorem~\ref{thm:gaps} each side contributes
$\sum_{j=1}^{\lfloor (h-1)/2\rfloor}j^{2}$. For $h$ odd this is
$2\cdot\frac{(\frac{h-1}2)(\frac{h+1}2)h}{6}=\frac{h(h^{2}-1)}{12}$; for $h$ even
it is $2\cdot\frac{(\frac h2-1)(\frac h2)(h-1)}{6}=\frac{h(h-1)(h-2)}{12}$.
\end{proof}

Theorem~\ref{thm:gaps} and Lemma~\ref{lem:w2} are exact and unconditional. They
are the reason the exponent $3/2$ appears below: it is forced by the triangular
weights, not fitted.

\begin{remark}
Theorem~\ref{thm:gaps} does \emph{not} evade Proposition~\ref{prop:info}: the
window position $m(n)$ and its length $h$ are determined by
$(\Mert(n),\Sq(n))$, and the primes are deterministic, so $E_{\mathrm{mid}}$
remains a function of $(\Mert(n),\Sq(n))$. What the theorem changes is what the
\emph{size} of $E_{\mathrm{mid}}$ is controlled by, namely the local second-order
statistics of prime gaps.
\end{remark}

\section{Two variance models}

Theorem~\ref{thm:gaps} suggests modelling $E_{\mathrm{mid}}$ as a weighted sum
of gap fluctuations. Two models are compared; both are heuristics, and both were
tested against the data of Section~\ref{sec:numeric}.

\begin{definition}[Independent-gap model]
Treat the $g_i$ near index $m(n)$ as uncorrelated with common local variance
$\sigma_g(n)^{2}$. Then by Lemma~\ref{lem:w2}
\[
\operatorname{SD}_{\mathrm{ind}}(n):=\sigma_g(n)\sqrt{\textstyle\sum_i w_i^{2}}
\;\approx\;\frac{\sigma_g(n)}{\sqrt{12}}\,h^{3/2}.
\]
\end{definition}

\begin{definition}[Covariance-corrected model]
Let $\rho(d)$ be the local autocorrelation of the gap sequence, and put
$R_{ab}=\rho(|i_a-i_b|)$ on the support of $w$. Define
\[
\kappa(h):=\frac{w^{\!\top}Rw}{w^{\!\top}w},
\qquad
\operatorname{SD}_{\mathrm{cov}}(n):=\sigma_g(n)\sqrt{\kappa(h)\textstyle\sum_i w_i^{2}} .
\]
\end{definition}

Because $w$ is antisymmetric about the centre, a \emph{positive} correlation
between gaps on opposite sides of the centre enters $w^{\!\top}Rw$ with a
negative sign. One therefore expects $\kappa(h)<1$, i.e.\ the independent model
to \emph{overpredict}. This is what is observed, and it accounts quantitatively
for the discrepancy (Section~\ref{sec:numeric}, part D).

\section{Numerical verification and measured laws}\label{sec:numeric}

\subsection*{Protocol}
$\mob$ was computed by a linear sieve for $n\le N=10^{6}$; $M,Q$ by prefix sums;
a prime table to index $304\,071$ ($p=4\,318\,813$); $T$ by prefix sums of the
primes. At $N=10^{6}$: $\Mert(N)=212$, $\Sq(N)=607\,926$, maximum required prime
index $304\,069$. All scripts and raw outputs accompany this note.

\subsection*{A. Exact statements verified}
The following hold with no exceptions for all $1\le n\le10^{6}$:
Lemma~\ref{lem:parity}; Proposition~\ref{prop:closed} (direct summation of $S$
against the closed form); Proposition~\ref{prop:block} and both inequalities of
Corollary~\ref{cor:bounds}; $\operatorname{sgn}T=\operatorname{sgn}\Mert$ and
$T=0\iff \Mert=0$ (there are $5\,361$ values of $n\le10^{6}$ with $\Mert(n)=0$);
Proposition~\ref{prop:offset}. Theorem~\ref{thm:gaps} was checked on $4\,000$
random $n\in[10^{3},10^{6}]$ with maximum absolute error $0$, and
Lemma~\ref{lem:w2} for all $h\le60$ and $h\in\{100,101,300,301,370\}$.

\subsection*{B. Global residual scale, $100\le n\le10^{6}$}
\[
\begin{array}{lcc}
\toprule
& E & E_{\mathrm{mid}}\\
\midrule
\text{median }|\cdot| & 1.09\times10^{3} & 9.74\times10^{2}\\
\text{r.m.s.} & 5.295\times10^{3} & 5.066\times10^{3}\\
\max|\cdot| & 6.2766\times10^{4} & 5.7654\times10^{4}\\
\text{median }|\cdot|/|T| & 1.069\times10^{-5} & 8.894\times10^{-6}\\
\bottomrule
\end{array}
\]

\begin{figure}[htbp]
\centering
\includegraphics[width=.85\textwidth]{figures/fig1_Emid.png}
\caption{$E_{\mathrm{mid}}(n)$ for $n\le10^{6}$.}
\end{figure}

\subsection*{C. The $h^{3/2}$ law and its stationarity}
Put $Z_{3/2}(n):=E_{\mathrm{mid}}(n)/\bigl(|\Mert(n)|^{3/2}\log n\bigr)$.
Conditioning on $h=|\Mert(n)|$ over $n\ge10^{3}$:
\[
\begin{array}{lcccccccc}
\toprule
h\in & [8,12) & [12,19) & [19,31) & [31,51) & [51,84) & [84,138) & [138,225) & [225,400)\\
\midrule
\operatorname{median}|Z_{3/2}| & 0.1426 & 0.1477 & 0.1589 & 0.1589 & 0.1556 & 0.1546 & 0.1507 & 0.1550\\
\text{count} & 36\,451 & 58\,125 & 86\,041 & 127\,696 & 152\,059 & 164\,880 & 231\,986 & 65\,791\\
\bottomrule
\end{array}
\]
and conditioning on $n$ instead:
\[
\begin{array}{lccc}
\toprule
n\in & \operatorname{RMS}Z_{3/2} & \operatorname{median}|Z_{3/2}| & \text{count}\\
\midrule
{}[1.00,1.58)\times10^{5} & 0.2331 & 0.1549 & 51\,624\\
{}[1.58,2.51)\times10^{5} & 0.2379 & 0.1574 & 82\,660\\
{}[2.51,3.98)\times10^{5} & 0.2340 & 0.1541 & 134\,286\\
{}[3.98,6.31)\times10^{5} & 0.2314 & 0.1530 & 215\,541\\
{}[6.31,10.0)\times10^{5} & 0.2324 & 0.1552 & 359\,873\\
\bottomrule
\end{array}
\]
Both tables are flat: the normalisation $|\Mert|^{3/2}\log n$ is stable across a
factor of $46$ in $h$ and a factor of $10$ in $n$. This supports, but does not
prove, the hypothesis that $Z_{3/2}$ converges in distribution.

\begin{figure}[htbp]
\centering
\includegraphics[width=.48\textwidth]{figures/fig2_scaling.png}
\hfill
\includegraphics[width=.48\textwidth]{figures/fig3_stationarity.png}
\caption{Left: $|E_{\mathrm{mid}}|$ against $|\Mert|$ with the envelope
$0.155\,h^{3/2}\log n$. Right: stationarity of $Z_{3/2}$ in $n$.}
\end{figure}

\subsection*{D. Which variance model is right}
Local gap statistics were measured on the prime-index band
$[5\times10^{4},3.05\times10^{5})$, which covers the windows for
$n\in[10^{5},10^{6}]$: mean gap $14.5927$, $\log p=14.6984$ at the band centre,
$\operatorname{Var}(g)=145.93$, hence
\[
\operatorname{Var}(g)/\log^{2}p=0.6755
\]
against the value $1$ predicted by a Poisson (Cram\'er) gap model. The gap
autocorrelation is negative at short lag and positive over a long tail:
\[
\begin{array}{lcccccc}
\toprule
d & 1 & 2 & 3 & 4 & 5 & 6\\
\midrule
\rho(d) & -0.0463 & -0.0225 & -0.0132 & -0.0057 & -0.0069 & -0.0061\\
\bottomrule
\end{array}
\]
with $\sum_{1\le d\le420}\rho(d)=+0.4646$. The resulting correction factor
decreases steadily,
\[
\begin{array}{lccccccc}
\toprule
h & 5 & 10 & 20 & 40 & 80 & 160 & 320\\
\midrule
\kappa(h) & 1.0206 & 0.9532 & 0.9094 & 0.8514 & 0.7869 & 0.7381 & 0.6871\\
\bottomrule
\end{array}
\]

Testing both models on the $883\,123$ values of $n$ with $n\ge10^{4}$ and
$h\ge12$:
\[
\begin{array}{lcccc}
\toprule
\text{model} & \operatorname{RMS}(E_{\mathrm{mid}})/\operatorname{RMS}(\mathrm{SD}) & \operatorname{RMS}Z & \operatorname{mean}Z & \operatorname{median}|Z| \\
\midrule
\text{independent gaps} & 0.8527 & 0.8963 & -1.88\times10^{-4} & 0.5962\\
\text{covariance-corrected} & 1.0031 & 1.0036 & -2.13\times10^{-4} & 0.6694\\
\bottomrule
\end{array}
\]
The $15\%$ deficit of the independent model is therefore explained, to within
$0.3\%$, by the correlations between prime gaps; and the corrected model
reproduces the Gaussian value $\operatorname{median}|Z|=0.6745$ to within $0.8\%$
with no free parameter. This settles the question left open by the independent
model.

\begin{figure}[htbp]
\centering
\includegraphics[width=.48\textwidth]{figures/fig4_kappa.png}
\hfill
\includegraphics[width=.48\textwidth]{figures/fig5_qq.png}
\caption{Left: the correction factor $\kappa(h)$. Right: quantiles of
$Z=E_{\mathrm{mid}}/\mathrm{SD}$ against the normal, on a thinned sample; the
independent-gap curve is visibly compressed, the corrected one is not.}
\end{figure}

\subsection*{E. Distribution of the standardised residual}
With $Z:=E_{\mathrm{mid}}/\operatorname{SD}_{\mathrm{cov}}$, by blocks in $n$
($h\ge12$):
\[
\begin{array}{lcccccc}
\toprule
n\in & \text{count} & \operatorname{mean} & \operatorname{sd} & \text{skew} & \text{exc.\ kurt.} & \operatorname{median}|Z|\\
\midrule
{}[10^{4},10^{5}) & 67\,120 & +0.0043 & 0.9813 & -0.0098 & +0.2556 & 0.6474\\
{}[10^{5},3.16\times10^{5}) & 183\,261 & -0.0000 & 1.0029 & -0.0167 & +0.1559 & 0.6705\\
{}[3.16\times10^{5},10^{6}) & 632\,741 & -0.0008 & 1.0061 & -0.0188 & +0.0957 & 0.6715\\
\bottomrule
\end{array}
\]
and by blocks in $h$, the excess kurtosis decreases monotonically,
\[
\begin{array}{lccccccc}
\toprule
h\in & [12,19) & [19,31) & [31,51) & [51,84) & [84,138) & [138,225) & [225,400)\\
\midrule
\text{exc.\ kurt.} & +0.676 & +0.338 & +0.213 & +0.031 & +0.058 & +0.024 & -0.143\\
\bottomrule
\end{array}
\]
which is the signature expected if a central limit theorem governs the sum of
Theorem~\ref{thm:gaps} as $h\to\infty$. Quantiles of $Z$ over all $n\ge10^{4}$,
$h\ge12$:
\[
\begin{array}{lccccccc}
\toprule
\text{level} & 1\% & 5\% & 25\% & 50\% & 75\% & 95\% & 99\%\\
\midrule
\text{observed} & -2.372 & -1.652 & -0.666 & +0.000 & +0.673 & +1.641 & +2.346\\
N(0,1) & -2.326 & -1.645 & -0.674 & 0 & +0.674 & +1.645 & +2.326\\
\bottomrule
\end{array}
\]

\subsection*{F. The points are not independent}
This is a necessary caveat on the preceding tables. Consecutive $n$ share almost
the same window, so the $883\,123$ values of $Z$ are strongly dependent: the
autocorrelation of $Z$ in $n$ is $+0.979$ at lag $1$, $+0.692$ at lag $10$, and
first falls below $0.05$ at lag $128$. The integrated autocorrelation time is
$\approx51$, giving an effective sample size of about $1.7\times10^{4}$, not
$8.8\times10^{5}$.

Consequently a Kolmogorov--Smirnov test on the full sample is meaningless. On
thinned subsamples:
\[
\begin{array}{lccccc}
\toprule
\text{thinning} & \text{count} & \operatorname{sd} & \text{exc.\ kurt.} & \mathrm{KS} & p\\
\midrule
1 & 883\,123 & 1.0036 & +0.120 & 0.00317 & 3.8\times10^{-8}\\
51 & 17\,317 & 1.0036 & +0.170 & 0.00451 & 0.87\\
128 & 6\,900 & 1.0057 & +0.018 & 0.00775 & 0.80\\
\bottomrule
\end{array}
\]
At the effective sample size, normality of $Z$ is not rejected. The apparent
rejection on the full sample is an artefact of dependence, not evidence against
the model.

\begin{figure}[htbp]
\centering
\includegraphics[width=.55\textwidth]{figures/fig6_autocorr.png}
\caption{Autocorrelation of $Z$ in $n$; the effective sample is roughly $50$
times smaller than the nominal one.}
\end{figure}

\subsection*{G. Relative size of the residual}
Combining $|E_{\mathrm{mid}}|\approx c\,h^{3/2}\log n$ with $|T|\approx
h\,p^{*}\approx \tfrac{3}{\pi^{2}}h\,n\log n$ predicts
\[
\frac{|E_{\mathrm{mid}}(n)|}{|T(n)|}\approx \frac{c\pi^{2}}{3}\cdot\frac{\sqrt{|\Mert(n)|}}{n},
\qquad \frac{c\pi^{2}}{3}=0.510 \text{ for } c=0.155 .
\]
Measured medians of $\bigl(|E_{\mathrm{mid}}|/|T|\bigr)\,n/\sqrt{|\Mert|}$:
$0.4385,\ 0.4800,\ 0.4995,\ 0.4955$ on
$[10^{3},10^{4}),[10^{4},10^{5}),[10^{5},3.16\times10^{5}),[3.16\times10^{5},10^{6})$.
This quantitatively explains why the central term reproduces $T$ so accurately,
and why $\operatorname{median}|E_{\mathrm{mid}}|/|T|=8.9\times10^{-6}$: if
$|\Mert(n)|$ is of size roughly $n^{1/2}$ then the relative correction is of size
$n^{-3/4}$.

\subsection*{H. Largest excursions}
The extremes of $E_{\mathrm{mid}}$ on $n\le10^{6}$ occur in tight clusters,
e.g.\ $-5.7654\times10^{4}$ at $n=924\,067$ ($\Mert=-301$) and
$+4.6576\times10^{4}$ at $n=932\,737$ ($\Mert=-269$). These are the $n$ at which
$|\Mert|$ is locally largest, consistent with the $h^{3/2}$ law and carrying no
separate information.

\section{Status of the claims}

To keep the register unambiguous:
\begin{itemize}[leftmargin=1.6em]
\item \emph{Proved}: Lemmas~\ref{lem:counters}, \ref{lem:parity}, \ref{lem:w2};
Propositions~\ref{prop:closed}, \ref{prop:block}, \ref{prop:scale},
\ref{prop:info}, \ref{prop:offset}; Corollary~\ref{cor:bounds};
Theorems~\ref{thm:RH}, \ref{thm:gaps}.
\item \emph{Measured}, over $10^{4}\le n\le10^{6}$ with effective sample
$\approx1.7\times10^{4}$: the law
$|E_{\mathrm{mid}}|\approx0.155\,|\Mert|^{3/2}\log n$; the failure of the
independent-gap model by a factor $0.853$ and its repair by the measured gap
covariance; approximate normality of the standardised residual, with excess
kurtosis decaying in $h$.
\item \emph{Conjectural}: that $Z_{3/2}$ converges in distribution; that the
limit is Gaussian; that a central limit theorem holds for the antisymmetric
weighted gap sum of Theorem~\ref{thm:gaps}. None of these is proved here, and
none would follow from the numerics alone.
\end{itemize}

\section{What the construction does and does not deliver}

It delivers an exact, elementary dictionary between $SS$ and $(\Mert,\Sq)$, a
two-sided transfer of the Mertens criterion (Theorem~\ref{thm:RH}), and a
residual whose size obeys a clean measured scaling law.

It does not deliver an independent arithmetic constraint on $\Mert(n)$, by
Proposition~\ref{prop:info}. For a transformed representation to become useful
it would have to produce at least one of:
\begin{enumerate}[label=(\roman*)]
\item a cancellation unavailable in the original variables;
\item a positivity or monotonicity principle strong enough to yield new bounds;
\item a spectral representation turning a Mertens bound into an operator inequality;
\item an inequality between independently controlled components;
\item a mechanism by which the prime side constrains the M\"obius side beyond
what $(\Mert,\Sq)$ already encodes.
\end{enumerate}
None of these has been extracted here. The measured law of
Section~\ref{sec:numeric} is of type (iv) in form only: it bounds
$E_{\mathrm{mid}}$ in terms of $|\Mert|$, which is the wrong direction to be
useful.

\section{A variant that is not an endpoint statistic (speculative)}

The obstruction identified in Proposition~\ref{prop:info} is caused by the use
of two \emph{separate} ordinal counters, which destroys the ordering of the
$\pm1$ occurrences. The minimal modification that removes it is a
\emph{shared} counter:
\[
\widetilde S(k):=\mob(k)\,p_{\Sq(k)},
\qquad
\widetilde{SS}(n):=\sum_{k\le n}\mob(k)\,p_{\Sq(k)} .
\]
Here the prime index depends on $k$ rather than on the sign-specific ordinal, so
$\widetilde{SS}(n)$ is not a function of $(\Mert(n),\Sq(n))$: by Abel summation
with $w(k)=p_{\Sq(k)}\asymp k\log k$,
\[
\widetilde{SS}(n)=\Mert(n)w(n)-\sum_{k<n}\Mert(k)\bigl(w(k+1)-w(k)\bigr),
\]
so it depends on the whole trajectory of $\Mert$ and not only on its endpoint.
This is offered as a direction, not a result; it is close to the classical
weighted Mertens sums $\sum_{k\le n}\mob(k)k^{\alpha}$ and may well reduce to
them.

\section{Novelty status}

The three zeta ingredients are classical. A search did not locate the sequence
$S$ ($2,-2,-3,1,-5,3,-7,1,1,5,-11,1,\dots$) as a named sequence or standard
transform. Absence of a search result is weak evidence and the search was not
exhaustive; this is recorded as an observation about the formulation, not a
claim about its consequences, which by Proposition~\ref{prop:info} follow from
$(\Mert,\Sq)$.

\section{Summary}

\[
\mob \longrightarrow (N_-,N_0,N_+) \longrightarrow S \longrightarrow SS,
\qquad
N_\pm=\frac{\Sq\pm\Mert}2,\quad N_0=n-\Sq,\quad \Sq\equiv\Mert\ (2),
\]
\[
\boxed{\;
SS(n)=\PS\!\left(\frac{\Sq+\Mert}2\right)-\PS\!\left(\frac{\Sq-\Mert}2\right)+n-\Sq \;}
\]
\[
\boxed{\;T(n)=SS(n)-n+\Sq(n)=\operatorname{sgn}(\Mert)\!\!\sum_{j\in W(n)}\!\!p_j,
\quad |W(n)|=|\Mert(n)|,\ \text{centre } \tfrac{\Sq(n)+1}2 \;}
\]
\[
\boxed{\;T(n)=\Mert(n)\,p_{\lfloor \Sq(n)/2\rfloor}\bigl(1+\theta(n)\bigr),\quad \theta(n)\to0\;}
\]
\[
\boxed{\;\mathrm{RH}\iff T(n)=O_\varepsilon\bigl(n^{3/2+\varepsilon}\bigr)\ \ \forall\varepsilon>0\;}
\]
\[
\boxed{\;T(n)=\Mert(n)\,p^{*}(n)+E_{\mathrm{mid}}(n),
\qquad
E_{\mathrm{mid}}(n)=\operatorname{sgn}(\Mert)\sum_i w_i g_i,
\qquad
\sum_i w_i^2=\frac{h^3}{12}\bigl(1+O(h^{-1})\bigr)\;}
\]
\[
\boxed{\;\operatorname{SD}\bigl(E_{\mathrm{mid}}\bigr)\;=\;\sigma_g(n)\sqrt{w^{\!\top}\!R\,w}\,,
\qquad
|E_{\mathrm{mid}}(n)|\approx0.155\,|\Mert(n)|^{3/2}\log n\;}
\]
the first line exact, the second measured over $10^{4}\le n\le10^{6}$ and not proved.

Every quantity above is a function of $(\Mert(n),\Sq(n))$. The transform is an
exact, prime-amplified reparametrisation of the Mertens function; the residual
$E_{\mathrm{mid}}$ trades the arithmetic of $\mob$ for the statistics of prime
gaps in a window of length $|\Mert(n)|$, at the cost of a factor
$|\Mert(n)|^{3/2}$, and therefore cannot improve a bound on $\Mert(n)$.

\end{document}
